\documentclass[11pt]{amsart}
\usepackage[margin=1in]{geometry}
\usepackage{amsmath,amssymb,array,booktabs,longtable}
\usepackage{xurl}
\usepackage[hidelinks,hypertexnames=false]{hyperref}
\newcommand{\GC}{\operatorname{GC}}
\newcommand{\Z}{\mathbb Z}
\newcommand{\record}[1]{\nolinkurl{#1}}
\newtheorem{theorem}{Theorem}
\newtheorem{corollary}[theorem]{Corollary}
\newcommand{\claimid}[1]{\par\nopagebreak\noindent\textit{Machine-readable claim:} \record{#1}.\par}
\setlength{\emergencystretch}{3em}
\raggedbottom
\title[Two-stage classification of linear gSW paths]{A two-stage classification of linear generalized Stone--Wales paths in Goldberg--Coxeter images of the dodecahedron}
\author{[human author]}
\date{Internal draft v6.2, 19 October 2026}
\begin{document}
\begin{abstract}
We study literal linear generalized Stone--Wales (gSW) paths in the duals
of $\GC(k,l)[C_{20}]$, where $k>l>0$. Such a path has distinct vertices,
an even number of vertices, degree-five endpoints, degree-six neighbours
of its endpoints, and triangular consecutive triples. Internal
degree-five vertices are not excluded by the definition. We present a
proof draft of an explicit two-stage classification. Nonprimitive
parameters give no path. For primitive parameters, put $m=\min(l,k-l)$.
Apart from unit-width cases, existence is equivalent to $k\equiv\pm1
\pmod m$ when $\lfloor k/m\rfloor\ge3$, and to $m\equiv\pm1\pmod r$,
where $r=k-2m$, when that quotient is two. Eight periodic belts give
positive witnesses; collisions preceding any admissible endpoint rule
out the other cases. Consequences are a smallest primitive negative
example $\GC(17,5)[C_{20}]=C_{7980}$, asymptotic density one of negative
primitive parameter points, and eventual periodicity along integer
affine rays. Forty-one arithmetic lemmas have been checked in Lean. A bounded
prototype additionally proves a literal path in an explicit unit-width
chart model. The general GC surface and complete classification remain
unformalized and retain prose-review obligations. An appendix records a
separate finite census of IPR fullerenes from $C_{60}$ to $C_{200}$:
120,622,224 outputs from each of two generators, with 816 negative
isomorphism classes. We distinguish finite checks, AI prose review,
formal proofs, and independently constructed verification. Novelty
relative to the full position-mapping and intersection theory remains
unsettled.
\end{abstract}
\maketitle

\noindent\textbf{Editorial status.} This is an unpublished, AI-assisted
proof draft, correcting the review metadata of English v6 and continuing Japanese v4. The theorem statements below
retain v4's numbering; calling a statement a theorem does not assert
completed external peer review. The coordinator (Claude) reviewed the
Japanese v4 text of Appendices B and C at 03:46 JST on 19 October 2026
(commit \record{49430e1c}), recalculating its formulas and recording no
missing steps. This English text transcribes that Japanese argument;
selected English formulas were also machine-checked by the coordinator.
This is AI prose review, not human peer review or a full formal proof.
The coordinator's later addendum also completes AI prose review of A-sheet items 2 and 3; Section 11 records the scope. The author placeholder and the AI-use statement are
for the human author's decision. The separate treatment of other
families and the general ray theorem is reserved for a follow-up paper.

\section{Objects and definitions}\label{sec:definitions}
\subsection{The literal path definition}
A fullerene is a simple, three-connected, cubic plane graph all of whose
faces are pentagons or hexagons. Its dual is a spherical triangulation
with twelve degree-five vertices and all other vertices of degree six.
Following Definition 1(a) of \cite{BBS}, a literal gSW path is a sequence
$p_0,\ldots,p_e$ satisfying:
\begin{enumerate}
\item $e\ge3$ is odd and all vertices are distinct;
\item $p_i$ is adjacent to $p_{i+1}$ and to $p_{i+2}$ whenever defined;
\item $p_0,p_e$ have degree five, and $p_1,p_{e-1}$ have degree six.
\end{enumerate}
The source's notation $v_1,\ldots,v_{2w}$, $w\ge2$, gives $e=2w-1$.
Consecutive adjacency is part of the word ``path''; distance-two
adjacency is explicit. In particular, internal degree-five vertices
are allowed. Requiring an even number of vertices is a definition, not
a new necessary condition. ``Positive'' and ``negative'' mean presence
and absence of this path. Reaction barriers, radial gSW operations,
and operations permitting heptagons are outside our scope.

\subsection{The A3 family and lattice coordinates}
Write $A3(k,l)=\GC(k,l)[C_{20}]$, $k>l>0$. Graver's older coordinates
$A3(p,q)$ correspond to $(k,l)=(p+q,p)$. Develop each parent face of
the icosahedral dual in the triangular lattice with
\[
 O=(0,0),\quad e=(k,l),\quad f=(-l,k+l),\quad T=k^2+kl+l^2.
\]
The fullerene has $20T$ vertices and its dual has $10T+2$. Let
$R(x,y)=(-y,x+y)$ and $G(u,v)=ue+vf$. The maps $G$ and $R$ commute,
and the inverse coordinate change is
\begin{equation}\label{eq:1}
 u=\frac{(k+l)x+ly}{T},\qquad v=\frac{-lx+ky}{T}.\tag{1}
\end{equation}
In parent coordinates write
\[
 U(i,j)=[(i,j),(i+1,j),(i,j+1)],\quad
 D(i,j)=[(i+1,j),(i+1,j+1),(i,j+1)].
\]
Developments use integer powers of $R$ and integer translations. A fine
unit triangle need not lie in one parent face: pieces cut by parent
edges are mapped separately. If $\gcd(k,l)=1$, the vectors $e,f,f-e$
are primitive, so no fine lattice vertex lies inside a parent edge.
Identifications of ordinary vertices can therefore be tested in face
interiors; parent vertices require separate treatment. Distinct face
labels do not imply disjoint closures.

\subsection{Kinds of evidence}
An AI proof draft is a general argument awaiting review. AI prose review
means the specified formulas or paragraphs were read and rederived by
Claude; it is neither human peer review nor necessarily an independent
derivation. Lean checks the stated formal propositions, separately from
the geometric hypotheses. Finite checks on an independent construction
are identified by both the constructor and the checker. ``R004
independent'' refers to completed reports from a verification team that
did not read the main team's proof or implementation within the stated
scope. The logical assembly of the classification has been checked in
prose, including the general cap, fan, cut, and ownership arguments
(P006-11). Sections 9 and 11 give the precise boundary.

\section{The classification}\label{sec:classification}
\begin{theorem}[Two-stage classification; proof draft]\label{thm:classification}
Let $k>l>0$ and $d=\gcd(k,l)$. If $d>1$, then $A3(k,l)$ is negative.
If $d=1$, put $m=\min(l,k-l)$. The graph is positive if and only if
one of the following holds:
\begin{equation}\label{eq:2}
\begin{aligned}
 &m=1;\\
 &m\ge2,\quad q=\lfloor k/m\rfloor\ge3,\quad k\equiv\pm1\pmod m;\\
 &m\ge2,\quad q=2,\quad r=k-2m,\quad
       [\,r=1\ \text{or}\ m\equiv\pm1\pmod r\,].
\end{aligned}\tag{2}
\end{equation}
\end{theorem}
\claimid{thm:classification}
For primitive parameters with $m\ge2$, equality $k=2m$ is impossible;
thus $0<r<m$ in the last line. No congruence modulo zero is used.
We handle $r=1$ before the remaining congruences.

The proof separates the integer-inflation obstruction (Section 3),
the two start orbits (Section 4), positive paths (Section 5), and
collisions with endpoint exclusion (Sections 6--7). The four regions
in the original coordinates are as follows.
\begin{center}\small
\begin{tabular}{lll}\toprule
Region & Positive construction & Negative construction\\\midrule
$k>3l$ & AR vertical/horizontal belts & Outer fan\\
$2l<k<3l$ & AT second belt & Positive fan, $c=2$\\
$l<k<2l$, $k\ge3(k-l)$ & Diagonal belt & Positive fan, $c=1$\\
$l<k<2l$, $k<3(k-l)$ & Folded second belt & Negative fan\\\bottomrule
\end{tabular}\end{center}
Primitive points with $k=2l$ or $3l$ have $l=1$ and belong to
Section 5.1. In the third region, $k\bmod(k-l)=l\bmod(k-l)$.
In the fourth, $r=2l-k$ and $m=k-l=l-r$, so $m\bmod r=l\bmod r$.
This gives (2), with the assembly reviewed in AU, Section 7.
The equality of decisions at $(k,l)$ and $(k,k-l)$ follows from the
classification, not from a graph isomorphism: their norms differ by
$2k(k-2l)$ and their orders generally differ.

\section{The obstruction for nonprimitive parameters}\label{sec:inflation}
\claimid{sec:inflation}
If $d=\gcd(k,l)>1$, the construction is the integer $d$-fold inflation
of $\GC(k/d,l/d)[C_{20}]$. More generally, in $\GC(d,0)[F]$ the
degree-five parent points lie in the $d$-multiple lattice in each chart.
Develop the path from its initial triangle. At an odd position, the
determinant of the lattice point with the direction of progression is
$\pm1$. If this point were degree five, that determinant would have to
be divisible by $d$, a contradiction. The gluings are lattice rotations
through multiples of $60$ degrees and transfers about points of the
$d$-multiple lattice, preserving divisibility. This does not assume a
global prohibition on internal degree-five vertices.

The arithmetic and transfer lemmas are formalized separately; Claude
reviewed the main integer-GC argument. R004 report 4 also gives an
independent general derivation. Neither constitutes a Lean definition
and proof of the entire GC surface. The original argument is
\record{general-inflation-argument.md}.

\section{Two starts and candidate endpoints}\label{sec:starts}
\claimid{sec:starts}
\subsection{Forced continuation and the first repeat}
Suppose every edge has exactly two common neighbours. From an oriented
triangle $(a,b,c)$, the next vertex is the common neighbour of $b,c$
other than $a$. Every simple path is a prefix of this forced walk. If
$\rho$ is its first repetition index, existence is equivalent to an
admissible endpoint before $\rho$, and $\rho\le |V|$. This conditional
principle is formalized as \record{R002.finite_literal_path_decision}.
Deriving the local assumption and the degree interpretation from the
surface, and covering all starts, remain outside that formal theorem.
The noncomputable \record{forcedWalk} is not a certified executable
classifier; see \record{formal-correspondence-decision.md}.

\subsection{The two rotational orbits}
The 60 orientation-preserving parent rotations act on the refined
lattice. Cyclic permutations of a face's corners act locally by
$R^{2p}z+(0,e,f)_p$ and respect the edge gluings. They are transitive
on the twelve parent vertices, with the fivefold stabilizer transitive
on the five initial radial edges. The two choices of third vertex have
opposite orientation and are not exchanged. Thus the 120 starts split
into two orbits of size 60, represented by
\[
 A=((0,0),(1,0),(0,1)),\qquad B=((0,0),(0,1),(1,0)).
\]
Their developments are
\begin{equation}\label{eq:3}
\begin{array}{ll}
 A: & P_{2j}=(0,j),\quad P_{2j+1}=(1,j),\\
 B: & P_{2j}=(j,0),\quad P_{2j+1}=(j,1).
\end{array}\tag{3}
\end{equation}
The AQ two-orbit argument was reviewed in prose. A negative proof must
exclude both representatives.

\subsection{The first parent-lattice candidates}
Assume $\gcd(k,l)=1$ and $l\ge2$. Since $k,l$ are coprime to $T$,
there is a unique $0<h<T$ with $kh\equiv l\pmod T$. The first odd
parent-lattice candidates in (1) are
\begin{equation}\label{eq:4}
 A:(1,h),\ E_A=2h+1;\qquad
 B:(t_B,1),\ t_B=T-h-1,\ E_B=2t_B+1.\tag{4}
\end{equation}
The first positive even candidate has index $2T$. The integer parent
coordinates $u_A,v_B$ satisfy
\begin{equation}\label{eq:5}
\begin{aligned}
 Tu_A&=k+l+lh,& Tv_B&=k-lt_B,\\
 ku_A&\equiv1\pmod l,& k(-v_B)&\equiv-1\pmod l.
\end{aligned}\tag{5}
\end{equation}
These formulas apply before a cone is encountered. They do not prove
that the candidates are reachable without repetition. Positive cases
require simplicity of the entire prefix; negative cases require a
separate exclusion of earlier parent vertices.

\section{Positive cases: eight belts and their caps}\label{sec:positive}
\claimid{sec:positive}
\subsection{The boundary \texorpdfstring{$l=1$}{l=1}}
Take the B sequence through $P_{2k+1}=(k,1)=e$. Its rectangle
$0\le x\le k$, $0\le y\le1$ lies in $U(0,0)\cup D(0,-1)$,
whose face labels differ. Each face map is injective and no fine vertex
lies inside the common parent edge. Only the endpoints are parent
vertices, giving an admissible path.

\subsection{Finite belt tables and infinite developments}
Let $F_z$ be the map from a developed parent triangle to its quotient
face. A period $v_0$ means $F_{z+v_0}(x)=F_z(x-v_0)$ on corresponding
cells: the affine map, not just the face label, repeats. Equality one
period apart and unique reversible edge gluing extend the table to all
integer cells.
\begin{center}\small
\begin{tabular}{llll}\toprule
Belt & Parent strip & $v_0$ & Faces per period\\\midrule
AR A & $0\le u\le1$ & $(0,5)$ &10\\
AR B & $-1\le v\le0$ & $(5,0)$ &10\\
AT A & $0\le v-2u\le1$ & $(2,4)$ &12\\
AT B & $0\le u+3v\le1$ & $(6,-2)$ &12\\
Diagonal A & $0\le v-u\le1$ & $(3,3)$ &12\\
Diagonal B & $0\le u+2v\le1$ & $(6,-3)$ &12\\
Folded second A & $-1\le v-2u\le0$ & $(2,4)$ &12\\
Folded second B & $-1\le u+3v\le0$ & $(6,-2)$ &12\\\bottomrule
\end{tabular}\end{center}
AR A alternates $U(0,j),D(0,j)$; AR B alternates
$D(j,-1),U(j+1,-1)$. A diagonal cell $i$ uses
$U(i,i),D(i,i),U(i,i+1),D(i,i+1)$. AT A uses $U(i,j),D(i,j)$ for
$j=2i,2i+1,2i+2$, and folded second A uses $j=2i-1,2i,2i+1$.
For each B belt rotate the triangle list by $R^{-1}$ and reglue from
its specified initial face; do not simply reuse rotated quotient labels.

The faces within a period are distinct, and the parent-vertex label
sets on the two boundaries are disjoint. Additional cap faces are
distinct from the central faces and from one another. There are 20
phase entries: $5+5$ for AR B, $1+2$ for AT A, $1+3$ for diagonal A,
and $1+2$ for folded second B. Their vertices, rotations, and translations
are recorded in \record{review-AW/cap-separation.md} and the fixed JSON
tables. T0743 checked the eight belts and T0746 the 20 cap entries on
Claude's independent construction. General containment is a separate step.

\subsection{Endpoints}
Use (3) up to index $E$ in the following table. Unit-width overlaps are
handled by a width-one or AR construction. If both signs coincide at
width two, choose the plus sign; the minus cap arguments use width at
least three.
\begin{center}\scriptsize
\begin{tabular}{llll}\toprule
Condition & Start & $Q$; endpoint index & Parent endpoint $W$\\\midrule
$l\ge2$, $k=ql+1$ & A & $q(k+l)+l$; $2Q+1$ & $(1,q)$\\
$l\ge2$, $k=ql-1$ & B & $q(k+l)+l$; $2Q-1$ & $(q+1,-1)$\\
$2l<k<3l$, $r=k-2l$, $l=sr+1$ & B & $s(3k+l)+k$; $2Q+1$ & $(3s+1,-s)$\\
Same region, $l=sr-1$ & A & $s(2k+3l)+k+l$; $2Q+1$ & $(s,2s+1)$\\
$l<k<2l$, $w=k-l$, $l=sw+1$ & B & $s(2k+l)+k$; $2Q+1$ & $(2s+1,-s)$\\
Same region, $l=sw-1$ & A & $s(k+2l)+k+l$; $2Q+1$ & $(s,s+1)$\\
$l<k<2l$, $w=2l-k$, $l=sw+1$ & A & $s(2k+3l)-k-l$; $2Q+1$ & $(s,2s-1)$\\
Same region, $l=sw-1$ & B & $s(3k+l)-k$; $2Q+1$ & $(3s-1,-s)$\\\bottomrule
\end{tabular}\end{center}
For AR, $kQ=qT+\varepsilon l$ and $lQ=T-\varepsilon(k+l)$.
For the remaining six rows, respectively,
\begin{equation}\label{eq:6}
\begin{aligned}
rQ+(3k+l)&=T,&rQ-(2k+3l)&=T,\\
wQ+(2k+l)&=T,&wQ-(k+2l)&=T,\\
wQ+(2k+3l)&=T,&wQ-(3k+l)&=T.
\end{aligned}\tag{6}
\end{equation}
In each A row $G(W)=(1,Q)$, and in each B row $G(W)=(Q,1)$,
except AR B, where it is $(Q-1,1)$. These integer identities are among
the Lean lemmas in Section 9; bounds, simplicity, and legality are not
consequences of the identities alone.

\subsection{Containment, including pieces cut by parent edges}
Use $[0,1]\times[0,Q]$ for A and $[0,Q]\times[0,1]$ for B, with
the upper horizontal bound $Q-1$ for AR B. These rectangles contain
all fine unit triangles used by the paths. Without caps, (1) and (6)
put the normal coordinate in the strip interval; for example, second
belt B has normal $[rx+(3k+l)y]/T\in[0,1]$.

For AR B, odd points above the strip have $t=0,\ldots,q-1$ and even
points below it have $t=Q-q,\ldots,Q-1$. The upper triangles have
vertices with $t\le q$. The bounds
\[
 T-k(q+1)-l=k(k-q-1)+kl+l^2-l>0,\quad
 k-q-1=q(l-1)-2\ge0
\]
put them in the initial face together with the belt. Point reflection
$P_{E-j}=P_E-P_j$ covers the terminal cap.
For second A, put $n=v-2u$. If $n<0$, then $x=ru-ln$ and $y\ge0$
give $0\le u\le1/r$, $-u\le v<2u$, and $u+v<3/r\le1$.
This lies in $U(0,0)\cup D(0,-1)$. Reflection $z\mapsto W-z$ sends
the $n>1$ cap back to this initial piece and gives terminal faces
$D(s-1,2s),U(s-1,2s+1)$.
For diagonal A, use $n=v-u$, $x=wu-ln$ to obtain
$0\le u\le1/w$, $-u\le v<u$, $u+v<2/w\le1$, and terminal faces
$D(s-1,s),U(s-1,s+1)$.

For folded second B, $n=u+3v$ and $y=ln-wv$. On the initial $n>0$
piece,
\[
 u=\frac{3y-(3l-w)n}{w}<\frac3w\le1,\qquad
 u+v=\frac{2y-kn}{w}<\frac2w<1.
\]
Both are nonnegative, so the piece lies in $U(0,0)\cup D(0,-1)$;
$D$ is central and $U$ additional. Reflection sends $n$ to $-1-n$.
These are planar containment arguments, not assertions of a reflection
symmetry of the quotient surface. Appendix C gives the terminal pieces
explicitly.

\subsection{Simplicity and degree conditions}
Two ordinary vertices identified in the central belt must have the
same face phase, so their developed difference is $aG(v_0)$,
$a\in\Z$. Across the width-one path rectangle the relevant transverse
period magnitudes are $5l,2r,3w,2w$ for AR, second, diagonal, and folded
second belts. Each exceeds one, excluding $a\ne0$; $a=0$ identifies
only the same index of (3). Cap labels are distinct and each face map
is injective, excluding further repeats.

Every construction has $Q<T$, excluding positive even parent points.
Substituting the possible integer normals $0,1$ (or $-1,0$) in (1)
and (6) leaves just the odd endpoint. For diagonal B with $w=1$, the
possible normal-zero point $(-2,1)$ is outside the rectangle. For
folded second A with $w=1$ it would have $u=-1$, again outside.
Minus-sign caps use width at least three and contain no other parent
point. The endpoints belong to disjoint boundary-label sets. Thus the
paths are simple, have even order, and meet all degree conditions.

\section{Negative cases: collisions after leaving a belt}\label{sec:negative}
\claimid{sec:negative}
The Japanese v4 B-side argument has coordinator AI prose review
(\record{49430e1c}). This English transcription retains the stated trust
limitations in Section 11.
The specified collision need not be the first repeat: an earlier
repeat is also an obstruction. We must follow every fine triangle up
to that time in the stated development and exclude an earlier legal
endpoint. Section 7 supplies the latter exclusion.

\subsection{The outer region}
Let $k=ql+r$, $q\ge3$, $2\le r\le l-2$, and put
\[
 a=(q-4)k+(q+2)l-1,\qquad b=(q+2)k+6l-1.
\]
A leaves $D(0,q)$ for $U(1,q),D(1,q),U(1,q+1),D(1,q+1)$;
the final two face maps compose the old belt map with
$H_A(z)=Rz+(q+2,-1)$. B leaves $U(q+1,-1)$ for
$D(q+1,-2),U(q+2,-2),D(q+2,-2),U(q+3,-2)$, using
$H_B(z)=Rz+(q+1,-q-2)$ on the final two. The rotation centres are
$(1,q+1)$ and $(q+2,-1)$. The finite table
\record{AT-A3-collision-atlas.json} covers $q\bmod5$.
The relations
\[
 H_A(q+5,z+(0,5))=H_A(q,z)+(0,5),\quad
 H_B(q+5,z+(5,0))=H_B(q,z)+(5,0)
\]
and belt periodicity extend it to all $q$.

The A rectangle $0\le x\le1$, $0\le y\le b$ has $0\le u<2$,
$v<q+2$. It crosses $u=1$ at $v=q+(r-x)/l\in(q,q+1)$.
For $u\le1$ one has $v<q+1$; for $u\ge1$, $v>q$. It is therefore
covered by the belt and the four added faces. The last point $(0,b)$
has $1<u<2$, $q+1<v<q+2$.
For B, $-2<v<1$ and $u+v<q+2$; its initial $v>0$ piece lies in
$U(0,0)$. At $v=-1$, $u=q+1+(r+y)/l\in(q+1,q+2)$, leading into
the four lower faces. The last point $(b,1)$ has $-2<v<-1$ and
$q+1<u+v<q+2$. Neither rectangle crosses the cut beyond the rotation
centre. The following six positive quantities justify these claims:
\begin{align*}
2T-k-l-lb&=(q^2-4)l^2+3qrl+2r^2-ql-r,\\
(q+2)T-kb&=(q-4)kl+(q+2)l^2+k,\\
(q+2)T-kb-k-l&=(q-4)kl+(q+2)l^2-l,\\
lb-T&=k(l-r)+5l^2-l,\\
kb-(q+1)T&=(4q-1)l^2+(q+5)rl+r^2-ql-r,\\
lb-T-k&=k(l-r-1)+5l^2-l.
\end{align*}
For $q=3$, substitute $k=3l+r$, $r\le l-2$ in the second and third
lines; for $q\ge4$, positivity follows directly. The collision identities are
\begin{equation}\label{eq:7}
\begin{aligned}
(1,a)&=R(0,b)+(q+2)e-6f,\\
(a,0)&=R(b,1)+(q-4)e-(q+2)f.
\end{aligned}\tag{7}
\end{equation}
These differ from the maps $H_A,H_B$ by the periods $-5f,-5e$.
The early A point has $0<u<1$ and the early B point $-1<v\le0$;
for example $T-k-l-la=kb-(q+1)T>0$. Since $0\le a<b$, the actual
revisits are $2b\to2a+1$ for A and $2b+1\to2a$ for B.

\subsection{The two positive-side fan regions}
Put $k=cl+w$, $l=sw+d$, $2\le d\le w-2$, with $c=1,2$ and
$s\ge2,1$, respectively. The strip $0\le v-cu\le1$ has period
$p(1,c)$, where $p=3,2$, respectively. Set
\[
\begin{split}
a&=(c+1)sk+(s-5)l-1,\qquad b=(cs+5)k+(c+1)sl-1,\\
W&=(s+1,cs+c+1),\qquad H(z)=R^{-1}z+(-W_v,W_u+W_v).
\end{split}
\]
From $D(s,W_v-1)$ pass through $U(s,W_v),D(s,W_v),U(s+1,W_v)$.
On the last two faces, the old belt map composed with $H$ is the new
map. B uses the rotated triangle list and translation, glued anew.
The table \record{AU-A3-fan-atlas.json} covers every $s\bmod p$.

Rotate B's fine and parent coordinates by $R$ only for comparison.
In common fine coordinates $(X,Y)$, A has $0\le X\le1$, $0\le Y\le b$;
B has $-1\le X\le0$, $0\le X+Y\le b$. Since
$X=wu-l(v-cu)$, the normal-one crossing is
$u=s+(d+X)/w\in(s,s+1)$, before the cut from $W$.
The bounds $v<W_v+1$, $u+v<s+W_v+2$ cover the portion beyond the
strip by the original cell and the three added triangles. If
$u\le s+1,v\le W_v$, use the original cell; if $v>W_v$, use the new
faces; if $u>s+1$, use the final $U$.
The initial negative-normal A piece has $u\le1/w$,
$-u\le v<cu$, $u+v<(c+1)/w\le3/4$, and lies in
$U(0,0)\cup D(0,-1)$. B has no such piece. The last points, A's
$(1,b)$ and B's original $(b,0)$, have $v>W_v$, $u+v>s+W_v+1$.

The nine quantities needed for these bounds, the early point, and
time ordering are strictly positive:
\begin{equation}\label{eq:8}
\begin{gathered}
a,\quad b-a,\quad T-b,\\
(cs+c+2)T-(kb+k+l),\\
((c+1)s+c+3)T-((k+l)b+k),\\
kb-l-(cs+c+1)T,\\
(k+l)b-((c+1)s+c+2)T,\\
T-[wa+(c+1)k+l],\qquad (s+2)T-(k+l+lb).
\end{gathered}\tag{8}
\end{equation}
Substitute $s=S+s_{\min}$, $w=D+E+4$, $d=D+2$, with
$S,D,E\ge0$, where $s_{\min}=2,1$. All 18 polynomials have
nonnegative coefficients and positive constant terms. Every permitted
parameter is obtained by $S=s-s_{\min}$, $D=d-2$, $E=w-d-2$.
This is a universal coefficient proof, not a finite extrapolation.
Lean checks the 18 inequalities, while their geometric interpretation
is separate. The collision identities are
\begin{equation}\label{eq:9}
\begin{aligned}
(0,a)&=R^{-1}(1,b)-(cs+5)e+(c+1)sf,\\
(a,1)&=R^{-1}(b,0)+(s-5)e+(cs+5)f.
\end{aligned}\tag{9}
\end{equation}
The images differ from $H$ by one belt period. The eighth quantity
puts the early point inside the old belt. Thus A revisits
$2b+1\to2a$, B revisits $2b\to2a+1$, and $0<a<b<T$.

\subsection{The negative-side second region}
Assume $3l/2<k<2l$, $w=2l-k$, $l=sw+d$, $s\ge2$,
$2\le d\le w-2$. The strip is $-1\le v-2u\le0$, with period
$(2,4)$. Set
\[
\begin{split}
a&=(2s-5)k+3sl-1,\qquad b=3sk+(s+5)l-1,\\
W&=(s+1,2s+1),\qquad H(z)=Rz+(3s+2,-s-1).
\end{split}
\]
Leave $D(s,2s)$ through $U(s+1,2s),D(s+1,2s),U(s+1,2s+1)$.
The last two maps use $H$. B uses the $R^{-1}$-rotated table.
\record{AU-A3-negative-fan-atlas.json} covers $s\bmod2$.
Now $X=-wu-l(v-2u)$, so at normal $-1$,
$u=s+(d-X)/w\in(s,s+1)$. The eight positive quantities are
\begin{equation}\label{eq:10}
\begin{gathered}
a,\quad b-a,\quad T-b,\quad (s+2)T-(k+l+lb),\\
(3s+3)T-((k+l)b+k),\quad lb-k-(s+1)T,\\
(k+l)b-(3s+2)T,\quad T-(wa+2k+3l).
\end{gathered}\tag{10}
\end{equation}
The same nonnegative substitution, with $s=S+2$, proves positivity
and is formalized by eight Lean lemmas. It gives $u<s+2$, $u+v<3s+3$
throughout the rectangle and $u>s+1$, $u+v>3s+2$ at the last point.
Outside the belt, $u\le s+1$ implies $v>2s-1$ and $u>s+1$ implies
$v>2s$, using
$k(s+1)-1-2sl=(s-1)w+2d-1>0$ in the latter case.
These bounds put every piece in the stated triangles without crossing
the cut. A has no initial positive-normal cap; B's initial cap lies
in the two faces of Section 5.4. Finally,
\begin{equation}\label{eq:11}
\begin{aligned}
(1,a)&=R(0,b)+3se-(s+5)f,\\
(a,0)&=R(b,1)+(2s-5)e-3sf.
\end{aligned}\tag{11}
\end{equation}
The difference from $H$ is one belt period, giving A's
$2b\to2a+1$ and B's $2b+1\to2a$ as actual ordinary-vertex repeats.

\section{Endpoint ordering and assembly}\label{sec:ordering}
\claimid{sec:ordering}
In the outer region, the nonunit residue and (5), with
$u_A\ge1$, $v_B<1$, imply $u_A\ge2$, $v_B\le-2$. Hence
\[
 lh=Tu_A-k-l\ge2T-k-l>lb,\qquad
 lt_B=k-Tv_B\ge k+2T>lb.
\]
Both odd candidates follow the collision. Also
$T-b=k(k-q-2)+l(k+l-6)+1>0$, excluding an earlier even parent point.

In the other regions, $k\equiv\pm w\pmod l$, $l=sw+d$.
For $1\le j\le s$, $w\le wj\le l-d\le l-2$; for $j=s+1$,
$wj=l+(w-d)$ with $2\le w-d\le w-2<l-1$. Thus none of
$1,\ldots,s+1$ gives $wj\equiv\pm1\pmod l$. Equation (5) yields
$u_A,-v_B\ge s+2$. The last bound of (8), or the fourth of (10),
again gives $h,t_B>b$, while $T-b>0$ excludes even parent points.
Apply this jointly with continuation of the development until a cone:
there is no other fine lattice point inside a unit triangle, and no
cone inside a parent edge. No hidden cone crossing is assumed away.
Both representative starts are excluded in every negative branch;
the two-orbit reduction excludes all 120 starts. Combining the positive
constructions, the region table, and the nonprimitive obstruction
proves the stated classification, subject to the review status in
Section 11.

\section{Three consequences}\label{sec:corollaries}
\subsection{The smallest primitive negative example}
\begin{corollary}\label{cor:minimum}
Among primitive parameters $k>l>0$, the smallest negative A3 example
is $(17,5)$, with $T=399$ and $20T=7980$.
\end{corollary}
\claimid{cor:minimum}
This is not minimality among all fullerenes or among nonprimitive A3
images. All 70 primitive pairs with $T<399$ satisfy (2); $(17,5)$
has quotient three and remainder two. Of the 123 pairs with $T\le673$,
the seven negatives are
\[
(17,5),(18,5),(19,7),(22,5),(17,12),(23,5),(21,8).
\]
The bound reduces enumeration to $k\le25$, $0<l<k$. Claude checked
all 123 on an independent construction (T0669); the other 116 were
positive. At $C_{7980}$ the two orbit repeats were $228\to15$ and
$229\to14$. See \record{A3-minimum-AP.md} and Section 10.3.

\subsection{Density of primitive parameter points}
\begin{samepage}\begin{corollary}\label{cor:density}
Let $D(X)$ count primitive pairs $k>l>0$ with $T\le X$, and let
$P(X),Z(X)$ count the positive and negative pairs. Then
\begin{equation}\label{eq:12}
\begin{aligned}
D(X)&=\frac{X}{\pi\sqrt3}+O(\sqrt X\log X),\\
P(X)&=\Theta(\sqrt X\log X),\\
\frac{Z(X)}{D(X)}&=1-\Theta(\log X/\sqrt X)\longrightarrow1.
\end{aligned}\tag{12}
\end{equation}
\end{corollary}
\claimid{cor:density}\end{samepage}
These are parameter counts, not a uniform measure on all fullerenes
or on graph isomorphism classes. The general counting argument remains
a proof draft, without complete independent review or formalization.
By the owner's decision of 19 October 2026, this density corollary is
excluded from the formalization target. Its counting, M\"obius-inversion,
error-bound, and limit machinery forms a separate substantial project;
the current priority is the graph classification, minimum, and affine-ray
corollaries. The prose proof and its unresolved review status are retained.
Put $K=\lfloor\sqrt X\rfloor$. Separate $m=1$, $k=qm\pm1$,
and the second-stage cases $r=1$ or $m=sr\pm1$,
$k=(2s+1)r\pm2$. There are at most two choices of $l$ for each
$(k,m)$, giving, with $H_K=\sum_{j=1}^K1/j$,
\[
 P(X)\le4K+(6K+8)H_K=O(\sqrt X\log X).
\]
For a lower bound, let $L=\lfloor\sqrt{X/3}\rfloor$ and
$M=\lfloor(L-1)/3\rfloor$. Count $(k,l)=(qm+1,m)$ with
$2\le m\le M$ and $3\le q\le\lfloor(L-1)/m\rfloor$.
These points are primitive and positive and satisfy $k\le L$,
$T\le3L^2\le X$. Their number is
$(L-1)(H_M-1)+O(M)=\Omega(\sqrt X\log X)$.
The region $k>l>0$, $T\le X$ has area $\pi X/(6\sqrt3)$,
with lattice-point error $O(\sqrt X+1)$. M\"obius inversion,
$\sum\mu(d)/d^2=6/\pi^2$, gives the formula for $D(X)$ with
error $O(\sqrt X\log X)$, and hence (12).

\subsection{Eventual periodicity on integer affine rays}
\begin{corollary}\label{cor:rays}
Let $k(t)=at+b$, $l(t)=ct+d$ have integer coefficients and satisfy
$k(t)>l(t)>0$ for all sufficiently large integer $t$. The existence
sequence is eventually periodic, with an arithmetically computable
threshold and period.
\end{corollary}
\claimid{cor:rays}
This is a consequence of Theorem 1, not a separate novelty claim.
If $\Delta=ad-bc\ne0$, then
$\gcd(k,l)=\gcd(k,l,\Delta)$, so primitivity depends only on
$t\bmod|\Delta|$. If $\Delta=0$, a nonconstant ray is an integer
multiple of a fixed primitive direction, with multiplier eventually
at least two, and is negative. A constant ray is constant.
Eventually $m=\min(l,k-l)$ is one affine expression. Constant $m$
gives fixed-modulus congruences. If $m$ grows, $\lfloor k/m\rfloor$
stabilizes; even when the limiting ratio is an integer, the constant
sign of the remainder selects the eventual quotient. At quotient two,
$r=k-2m$ is affine; persistent $r=0$ is nonprimitive and $r=1$ is
positive. The remaining conditions are affine divisibility tests.

For $A(t)=\alpha t+\beta>0$, $B(t)=\gamma t+\delta$ with
$\alpha>0$, $A\mid B$ implies $A\mid(\alpha\delta-\gamma\beta)$.
If that constant is nonzero, divisibility eventually fails. If it is
zero, $B=(\gamma/\alpha)A$ and divisibility holds identically exactly
when $\gamma/\alpha$ is an integer. For $\alpha=0$ it is a
fixed-modulus congruence. Combine the inequality thresholds and take
the least common multiple of the fixed moduli and $|\Delta|$.
This is the direct argument in \record{A3-ray-theorem-AU.md}; after
classification, the conclusion also falls under parametric Presburger
arithmetic \cite{BGW}.

\section{Correspondence with Lean (P001/P002)}\label{sec:lean}
\subsection{Formal inputs, conclusions, and missing connections}
The A3 arithmetic consists of 15 declarations in
\record{formal/Formal/A3Arithmetic.lean} and 26 in
\record{formal/Formal/A3PolynomialBounds.lean}, all in namespace
\record{R002.A3}. Variables are integers; \texttt{norm k l} is $T$,
\texttt{gc k l u v} is $G(u,v)$, and the formal statements call $T$
$N$. The rotations are \record{rotate60} and \record{invRotate}.
\begin{longtable}{>{\raggedright\arraybackslash}p{.30\textwidth}>{\raggedright\arraybackslash}p{.63\textwidth}}
\toprule Declaration(s) & Exact scope and what is not formalized\\\midrule\endhead
\record{unit_endpoint} & Arbitrary integers $q,l,\varepsilon$, not just
$\varepsilon=\pm1$; the two AR identities. Branch selection, endpoint
interpretation, simplicity, and degrees are separate.\\
\record{second_plus},\newline \record{second_minus} & Define $l=sr\pm1$, $k=2l+r$;
prove (6) and $G(W)$. The parameter region, small widths, caps, and
$Q<T$ are not conclusions.\\
\record{diagonal_plus},\newline \record{diagonal_minus} & Define $l=sw\pm1$,
$k=l+w$; corresponding identities and endpoints. Boundary labels and
absence of other cones are not formalized.\\
\record{folded_second_plus},\newline \record{folded_second_minus} & Define
$l=sw\pm1$, $k=2l-w$; corresponding identities. Negative normals,
cap ownership, and simplicity are outside the statements.\\
\record{outer_collisions},\newline \record{positive_fan_collisions},\newline
\record{negative_fan_collisions} & Identities (7), (9), (11), for arbitrary
integer arguments, including arbitrary $c$. No inequalities or
primitivity assumptions. Gluing, actual passage, and quotient equality
are additional obligations.\\
\record{no_small_multiple} & For $0<r<l$, excludes $r=lz$.\\
\record{outer_u_bound} & If $0<l$, $2\le r\le l-2$, $u\ge1$, and
$(ql+r)u-lv=1$, then $u\ge2$.\\
\record{outer_v_bound} & If $2\le l$, $2\le r\le l-2$, $v<1$, and
$lu+(ql+r+l)v=1$, then $v\le-2$. The exclusion through $s+1$ in
Section 7 is not covered by these three declarations.\\
\record{candidate_order} & Assumes $N,k\ge0$, $l>0$, $J\le u$,
$v\le-J$, $Nu=k+l+lh$, $N(-v)=lt-k$, and
$JN-k-l-lb>0$. Concludes $b<h$, $b<t$, $2b<2h+1$,
$2b+1<2t+1$. The geometric meaning and the bound on $J$ are inputs.\\
\record{outer_endpoint_order} & Combines the outer bounds with
$N=\operatorname{norm}(k,l)$, the coordinate equations (including
auxiliary $u_B,v_A$), $2N-k-l-lb>0$, and $b<N$; also concludes
$2b+1<2N$. The gap and geometric applicability remain assumptions.\\
\record{au_c1_*_pos}, \record{au_c2_*_pos} & Eighteen inequalities (8)
after $s=S+2$ or $S+1$, $w=D+E+4$, $d=D+2$, $l=sw+d$,
$k=cl+w$, with $S,D,E\ge0$. The inverse substitution and translation
to rectangle and cut containment are supplied in prose.\\
\record{au_neg_*_pos} & Eight inequalities (10), with $s=S+2$ and
$k=2l-w$. Ownership of the initial and final pieces is not formalized.\\
\bottomrule
\end{longtable}
The nine c1/c2 suffixes are
\record{a_pos}, \record{b_minus_a_pos}, \record{N_minus_b_pos},
\record{v_upper_pos}, \record{sum_upper_pos}, \record{v_late_lower_pos},
\record{sum_late_lower_pos}, \record{early_inside_pos},
\record{endpoint_bound_pos}. The eight negative suffixes replace the
two v-suffixes by \record{u_upper_pos}, \record{u_late_lower_pos} and
omit \record{endpoint_bound_pos}. Forty-one is a declaration count,
not a percentage of formalization. Theorem 1, the general GC surface, the
two-orbit reduction, and quotient simplicity are not Lean theorems.
The conditional finite-prefix theorem is separate from these lemmas.

\subsection{Environment, axioms, and saved execution}
The fixed environment is Lean 4.35.0 and mathlib v4.35.0, with manifest
commit \record{56a005855ea17e0b0958f73ec801ae140b0a1679}.
T0738 ran \texttt{env LEAN\_NUM\_THREADS=1 lake --no-cache build +Formal:olean}
from \record{formal}, exited zero, and used 33.786 seconds and peak
2836.5 MB with no charge. The base is
\record{76a9d89ebce23ad15975e2eaeb716ddd5f90fcfc} plus the trial's patch
and snapshots. The log says \emph{Replayed} for A3Arithmetic and
\emph{Built} for A3PolynomialBounds; this was not a fresh rebuild of
every declaration. Header, style, and unused-variable warnings remain.
The earlier technical failures T0736--T0737 are preserved.

Saved transitive \texttt{\#print axioms} output lists only
\record{propext} for \record{unit_endpoint}; \texttt{propext, Quot.sound}
for the other six endpoint and three collision declarations; and
\texttt{propext, Classical.choice, Quot.sound} for the five congruence
and ordering declarations and the 26 inequalities. No \record{sorryAx},
custom axiom, or additional native-evaluation axiom appears in those
complete lists. This is not an independent audit of the executable
or all caches. The source hashes are
\begin{quote}\small
A3Arithmetic:\par
\record{25a2259a80471a4000f5777a3cf5ef71c569d4ad00f8985f14ebe54ebbe41e57}\par
A3PolynomialBounds:\par
\record{e9a6ca756c53b72b05d0e2e05d83c2fe57a234ad4453df9cb88ea2c3dbd31407}\par
T0738 stdout:\par
\record{0489ab519ab3087a8816215d26eb92ca8e79bae5fd5798430d5fe31e54af3676}
\end{quote}
The v4 correspondence checked equality with the saved snapshots.
The v5 English rewrite did not rerun Lean or claim a separate-kernel audit.
The following v6 prototype has its own executions and scope.

\subsection{Unit-width prototype and current formalization target}
The next-56 prototype adds 904 Lean source lines (including its root
import), measured at 27 minutes 30 seconds of wall time in the existing
environment. It fixes twenty oriented parent faces, certifies reversed
edge gluing, complete five-corner links, and connectedness of the face
gluing. For $n=k-1$, its vertices are twelve parent labels and twenty
copies of the nonnegative triples $x+y+z=n$. Their coordinates $(x,y+1)$
are in bijection with the interior integer points of the parent triangle
$O=(0,0), A=(k,1), B=(-1,k+1)$. The parent edges have no other integer
points. Cartesian unit differences across the half-turn seam are proved
equivalent to the array adjacency, for all parameters.

\record{R002.GCUnit.unit_width_path} proves that for every $k\ge2$ this
explicit graph has a literal path with $2k+2$ vertices: prefix
injectivity, consecutive and skip adjacency, actual degrees five at
the ends and six next to the ends. These are conclusions of the graph
definition, not assumed cone labels.
\record{R002.GCUnit.Geometry.chartAdj_iff} establishes the Cartesian
adjacency correspondence; \record{R002.Icosahedron.faces_connected}
certifies the fixed face connection. This is a concrete unit-width
model theorem. Equivalence with an independent general GC surface
quotient, all ordinary degrees, the two-common-neighbours property,
and the complete classification remain separate obligations. The same
agent wrote both descriptions; this is not independent specification
review or a Comparator check.

T0924 built the integrated root and then re-elaborated the axiom-check
module without relying on its build cache. All 23 selected declarations
use only a subset of \texttt{propext, Classical.choice, Quot.sound};
the path and adjacency correspondence use those three. The build used
45.037 seconds and peak 2840.2 MB. T0926 confirmed rejection of the
existing unfinished-proof axiom. The command checks an allowlist of
three axioms, not a list of known bad names. No native evaluation or
added axiom is used in the prototype proofs. Failed attempts are
preserved. Style and module-migration warnings remain; no independent
kernel or package-release audit is claimed.

T0914 matched all vertices and edges of the old small-triangle
constructor for the already known $k=2,3,4$. This finite comparison
does not prove general GC equivalence. The command and exact scope are
in \record{formalization-unit-prototype.md}, and measurements in
\record{prototype-l1-measurements.json}. The current formalization target
is Theorem~\ref{thm:classification} and
Corollaries~\ref{cor:minimum} and~\ref{cor:rays}.
Corollary~\ref{cor:density} is excluded as stated in its section.
The revised estimate is \record{lean-full-formalization-estimate-v2.md}.

\section{Finite checks and selection history}\label{sec:checks}
\subsection{Predictions and comparisons}
Discovery data and later fixed predictions are distinguished; successful
later tests do not erase failed conjectures or selection history.
\begin{longtable}{p{.16\textwidth}p{.77\textwidth}}
\toprule Stage & Scope and result\\\midrule\endhead
AP/T0669 & All 123 primitive pairs through $T=673$, rechecked after
discovery on an independent construction: 116 positive, seven negative.
These are not 123 unseen tests.\\
AR & Twenty points fixed at \record{23486766}: 12 positive, eight negative,
all matching.\\
AS & The 921 points through $T=5000$ were used for discovery (602 positive,
319 negative). Fourteen separate predictions at \record{e379785c} all
matched. The three counterexamples to AQ's stronger threshold remain
in the history.\\
Claude T0708 & All 360 pairs with $5000<T\le7000$, independently
constructed: 171 positive, 189 negative. Twenty-six were already known,
so 334 were new.\\
AT & Forty-six points at \record{15e419a4}, $9049\le T\le185851$:
24 positive, 22 negative, all matching.\\
AU & Thirty points at \record{1ede36c9}, maximum $T=488721$: 12 positive,
18 negative; all decisions, 12 endpoints, and 36 negative-start
collisions matched.\\
AV & Cutoffs $10^4,10^5,10^6,10^7$ and the checking method fixed at
\record{bba3e303}, without predictions of exact finite counts. Two
arithmetic implementations agreed; no new graph census was run.\\
\bottomrule
\end{longtable}
At $X=10^7$, the arithmetic count was 1,837,762 primitive points,
81,946 positive and 1,755,816 negative (about 95.54\%). This is not
a proof of the limit or of monotonicity. Direct endpoint formulas
were compared with unit continuation on 161 AR, 76 AT, and 323 AU
paths. These 560 paths include repeated graphs.

\subsection{Belt tables and the AW package}
Claude independently recalculated the eight belt tables in T0743.
Execution of \record{review-AW} is recorded in AX: T0746 checked eight
belts and 20 caps; T0747 checked 337 endpoint orders (159 outer,
63 c2, 68 c1, 47 folded second); T0749--T0768 checked every vertex,
cut triangle piece, and cut for 160 known negative graphs and 320
starts. T0773--T0782 checked all 19,200 starts of those same graphs.
T0770 compared 1809 arithmetic values on nine rays; T0771 checked
14 graph instances. Probes T0748/T0769 are not additional examples.
T0772 stopped before execution on a resource-limit error. The 37 trials
comprise 36 successes and one technical failure: 178.167 seconds total,
maximum 13.817 seconds per trial and 97.3 MB, no charge. The aggregate
286,940 vertices, 286,300 triangles, 303,857 pieces, and 19,412 cuts
are workload counts, not independent examples. Codex wrote the checking
code with knowledge of the proof; Claude ran it on an independent
constructor. This is finite verification, not an independent general proof.

\subsection{Reports received from R004}
Earlier general integer-GC results use the completed reports 3--6 at
\record{234253238be3949656fa44d63829bbc3e16eb4a3}, fixed in AM. Reports
9--10 were received at \record{da213c7c}, with byte-identical copies in
\record{A3-v3-received}. For the census, report 9 established exact
agreement through $C_{176}$, with 19,703,439 outputs per generator and
228 negatives. R004 compared complete canonical-code sets, including
positive graphs. Its Colab checks reread every stored graph and its
decision evidence; Mac checks covered hashes, decision rows, full code
sets, and negative certificates, not a fresh check of every positive
path on the Mac. Report 10's partial $C_{178}$ recovery (34 of 256
partitions, 940,264 graphs, no negatives) was not a completed census.

\textbf{Update authorized for this draft.} The coordinator's next-51
instruction confirms independent agreement through $C_{182}$.
For this draft's fixed receipt boundary, $C_{184}$--$C_{200}$ remain
pending. Updated independent cumulative totals are deliberately not
inserted: they require the coordinator's final check. The main team's
fixed $C_{60}$--$C_{200}$ totals in Appendix A are unchanged. This
boundary is a dated receipt state, not a claim about live jobs.

R004 report 10 checked the A3 claim using a GC construction implemented
from the original source, both adjacency search and face-strip tracking.
For $T\le700$, the theorem domain $k>l>0$ contains 191 checked pairs:
127 primitive (120 positive, seven negative) and 64 nonprimitive
(all negative). A further 39 nonprimitive boundary pairs, 25 of the form
$(k,0)$ and 14 of the form $(k,k)$, are outside that domain and also
negative. Thus there are 230 non-control cases in total. The two controls
$(1,0),(1,1)$ give 232 stored records. The coordinator's v5 arithmetic
review corrected the earlier description of all 230 as theorem-domain
cases; see correction P006-09. All decisions matched, including
minimality at $T=399$. The report reread all records, positive paths,
and 13,200 negative-start certificates. The main team's proof and code
were not read, but the expected claim was known: this was not blind.
The same verifier used multiple implementations; it was not a further
agent's audit. Editing this paper did not rerun those constructions.

R004 independently reduced the nonprimitive case to integer GC
inflation. It also supplied a general positive proof for
$k\equiv\pm1\pmod l$ at \record{df84a3d2}, using a ten-face cylindrical
belt and two caps without consulting the main proof. The coordinator
rederived its identities, containment, distinctness, and degrees
(corrected receipt \record{0c59d19f}) and reported no gap. This is not
human review or formalization; independent necessity and the complete
primitive classification remain unfinished. Both proofs use belts and
transverse width, so methodological nonequivalence is not asserted.
Results on other families do not substitute for an independent A3 proof.

\subsection{Additional checks and the \texorpdfstring{$T=1099$}{T=1099} separation}
The \record{review-A3} package was fixed before execution at
\record{cb01852c}. Claude ran Codex's checks on its independent graphs.
T0789 verified 24 gluing records (10 outer, 10 positive-side, four
negative-side), 82 edge moves, 48 terminal triangles, and 144 vertex
images. T0795--T0864 checked 560 known positive paths on 417 distinct
graphs, including 584,420 vertices and 583,300 unit triangles.
T0865--T0884 checked 159 known outer-negative graphs, 318 starts,
including repeats, prior-cone exclusion, regions, cuts, and edge gluings.
The five probes T0790--T0794 overlap these samples. T0789--T0884
contains \emph{96}, not 95, successes: one atlas, five probes, and
90 batches. This corrected addition is recorded in
\record{A3-v2-validation.json}. Total time was 234.914 seconds (231.601
for the batches), maximum 7.034 seconds per trial and 103.8 MB, no
charge. T0788 was a pre-execution argument-expansion failure. None of
these known examples is counted again as an unseen test.

R005 compared $(23,15)$ and $(27,10)$: both have $T=1099$, the same
final group product $u$ on 60 states, and six closed zigzags of length
10990. Formula (2) nevertheless labels the first positive and the
second negative. This comparison initially used arithmetic, not graph
decisions. Claude subsequently constructed both $C_{21980}$ graphs
(10,992 dual vertices, IPR) independently, finding a 290-vertex path
in the first and no path in the second. It did not recalculate the
zigzag invariants; this rewrite did not rerun the graph tests.
The pair was selected during literature comparison, not preregistered
as an unseen sample. It shows that existence is not a function of
$(T,u)$ alone. It does not exclude a derivation using the full position
mapping, z-vectors, or intersection order.

\section{Review obligations and related work}\label{sec:review}
\subsection{Current status of the claims}
The definition, two start orbits, and candidate endpoints have been
checked against the source and in AI prose review, but lack a complete
formal surface connection. The main AR positive argument was reviewed;
its finite belt and cap tables were checked later. Through the fixed
review \record{9fedf3fc}, the coordinator checked congruence exclusion,
endpoint ordering, the positive-side A calculations, outer A coordinates,
negative-side A regions, the collision identities, and their period
differences. Ownership of all outer A pieces and some negative-side A
lower bounds retained abbreviated steps. The periodic transfer of the
finite tables was rederived. Signs, small widths, simplicity, and degrees
were checked, while some full cap-piece arguments remained partial.
Appendices B and C explicitly supply B's three regions and the terminal
caps. A later fixed review, \record{49430e1c} (03:46 JST, 19 October
2026), records the coordinator's recalculation of all formulas in the
Japanese v4 Appendices B and C, with no missing steps found. This review
was omitted from English v6's status metadata. English v6.2 corrects the
omission; it transcribes the Japanese proof, with the coordinator's
separate machine checks of selected English formulas recorded in P006.
The later addendum \record{c1b27f1d}, headed 15:10 JST, completes
the general AI prose review of A-sheet items 2 and 3: Appendix C
supplies the terminal containment, and the diagonal-A initial cap
was rederived in that addendum. P006-11 records this update.
The English-v5 review addendum \record{f88d1b6a} identifies its one
failed check as the R004 nonprimitive domain count corrected in
P006-09; the other 76 checks passed. The typesetting issue was
outside those 77 checks. Neither addendum supplies human peer
review, full formalization, or an independent necessity proof.

Thus the complete classification has finite agreement and 41 Lean
arithmetic lemmas, but its full geometric connection remains a proof
draft. R004's 230 independent finite cases and its general AR positive
proof do not close all primitive negative cases. The minimum has
independent finite confirmation; the general density and ray arithmetic
retain prose-review obligations. The census has the distinct receipt
boundary in Section 10.3. Human peer review, full formalization, and an
independent general proof of the complete classification have not occurred.

\subsection{Related work and the novelty boundary}
The problem comes from the classification question for fullerenes
admitting no generalized Stone--Wales operation in \cite{BBS}. Our
literal Definition 1(a) scope and the restriction to interior A3
parameters are essential. The Goldberg--Coxeter construction,
tightness of primitive images, and arithmetic of closed zigzags are
established theory. Dutour--Deza's position mapping uses a directed
parent-dual edge and a position $p$; its group product determines
zigzag number and length and can be computed by a Euclidean algorithm
\cite[Section 4, Theorems 4.5, 4.12, 6.5 and Proposition 4.14]{DD}.
For primitive dodecahedral images, the normalized partition vector has
three possible types. Zigzags, railroads, and their self- and mutual
intersections have also been studied \cite{ZZ,Knots}. None of these
constructions or methods is itself claimed as new here.

Closed-circuit invariants and the existence of a simple open gSW path
retain different information. In R005's comparison of 489 already
known primitive pairs, adding the parent-distance sequence of pentagon
contacts or the counts of shared pentagons still did not separate
the two-stage decisions. At $T=1099$, even the complete pentagon-contact
cycles of all zigzags, in the fixed parent labels and up to cyclic
start and reversal, agreed for the pair above. This does not assert
equality of all hexagon revisit orders or edge-intersection vectors.
The 489 labels were obtained from the AU formula, not a new census of
graphs. The report is \record{R005/report-2026-10-19-position-mapping.md}.

Full timing information can recover the input. Put $n=k+l$ and let
$\tau$ be the first position-update count from $p=1$ to $p=n$.
Then $\tau l\equiv-1\pmod n$, so the unordered adjacent intervals
$\{\tau,n-\tau\}$ determine $(k,l)$ under $k>l>0$.
Recovering the input alone is not a short deduction of gSW existence.
The additional geometric task is to preserve legal endpoints, parity,
and the order of revisits to original dual vertices, and reduce their
stopping condition to the two residue stages.

Comparing self-intersection against arrival at the next pentagon also
has precedent: \cite[Section 6, Fig.~5]{Growing} uses such a comparison
for Coxeter coordinates $(7,2)$, $C_{1340}$, to exclude a particular
twisted Endo--Kroto reduction patch. That path and operation differ from
literal gSW. Catalogues of isomerization patches \cite{Patches},
Caspar--Klug $T$- and gcd-based classifications \cite{Capsids}, and
congruence conditions in a pentagonal-polyhedron model \cite{Pentagonal}
are related, but the portions examined did not supply this same gSW
criterion. Reading of \cite{Patches,Knots} was limited to abstracts;
\cite{Pentagonal} was read mainly in its abstract and Section 2.
Unread material has not been ruled out.

The candidate contribution is the explicit two-stage necessary and
sufficient condition in the interior A3 domain, with geometric positive
witnesses and negative revisits. The literature search through
19 October 2026 did not locate the same criterion. Whether it is a
short consequence of the full position-mapping and intersection theory
remains open; no priority or worldwide novelty guarantee is made.
After classification, ray periodicity is a consequence of parametric
Presburger arithmetic \cite[Theorem 1.15, Property 1]{BGW}, not a new
general periodicity principle.

\appendix
\section{IPR census from \texorpdfstring{$C_{60}$}{C60} through \texorpdfstring{$C_{200}$}{C200}}\label{app:census}
\claimid{app:census}
\subsection{Scope, method, and preservation}
IPR means that no two pentagons share an edge. The 71 even sizes
$N=60,62,\ldots,200$ were enumerated with buckygen 1.1 and fullgen
from plantri 5.8. The saved cumulative record P contains 120,622,224
outputs per generator. The pipeline checked the dual, spherical
rotation system, degrees, IPR, and the classifier's local prerequisites.
Negative canonical-code sets agreed at every size, giving 816 distinct
classes. The main team did \emph{not} compare full canonical-code sets
of all positive graphs. Using two generators does not independently
validate their shared decision code.

Before adopting the C classifier, the C--Python audit compared it with
both forced continuation and a definition-based DFS on 1,456,598
general $C_{20}$--$C_{100}$ outputs and 4,770,156 IPR
$C_{60}$--$C_{160}$ outputs: 6,226,754 in total, all agreeing, with
positive witnesses checked against the definition. These were buckygen
outputs and the two ranges overlap. This is not a Python rerun of all
120,622,224 outputs. Subsequent negatives were checked by both Python
searches and exhaustive-start absence certificates. Generator
completeness and correctness on every possible input are not formally
proved by these tests.

Saved records include size/generator/partition counts and completion,
implementation hashes, negative full codes and adjacency lists, 16
audits and cumulative summaries, and one JSON certificate per negative
class. A certificate contains the dual adjacency and the maximal
simple forced path and stopping reason for every admissible initial
triangle. The index's \record{canonical_sha256} hashes canonical-code
bytes; \record{certificate_sha256} hashes the JSON file. The original
audits compared the full codes, not just their hashes. The v3 metadata
extraction verified saved hashes and counts without rerunning search,
canonicalization, or certificate logic.

The entire census was not one preregistered experiment: it includes
discovery and later confirmation. Predictions of the full negative
sets for $C_{186}$--$C_{200}$ were fixed at \record{3ac5c0b} and matched
unchanged. Failures, interruptions, and selections are preserved in
report P and its trials.

\subsection{Counts by size}
The following table transcribes the frozen \record{A3-v3-data/counts.csv}.
Zero total output and zero negatives among existing graphs are distinct.
% BEGIN GENERATED COUNTS: exact rows from the fixed CSV, not a new census.
\begin{longtable}{rrr}
\toprule Carbon order & Outputs per generator & Negative classes\\\midrule\endhead
60 & 1 & 0 \\
62 & 0 & 0 \\
64 & 0 & 0 \\
66 & 0 & 0 \\
68 & 0 & 0 \\
70 & 1 & 0 \\
72 & 1 & 0 \\
74 & 1 & 0 \\
76 & 2 & 0 \\
78 & 5 & 0 \\
80 & 7 & 1 \\
82 & 9 & 0 \\
84 & 24 & 0 \\
86 & 19 & 0 \\
88 & 35 & 0 \\
90 & 46 & 0 \\
92 & 86 & 0 \\
94 & 134 & 0 \\
96 & 187 & 1 \\
98 & 259 & 0 \\
100 & 450 & 0 \\
102 & 616 & 0 \\
104 & 823 & 1 \\
106 & 1,233 & 0 \\
108 & 1,799 & 0 \\
110 & 2,355 & 0 \\
112 & 3,342 & 2 \\
114 & 4,468 & 0 \\
116 & 6,063 & 0 \\
118 & 8,148 & 0 \\
120 & 10,774 & 3 \\
122 & 13,977 & 0 \\
124 & 18,769 & 0 \\
126 & 23,589 & 0 \\
128 & 30,683 & 6 \\
130 & 39,393 & 0 \\
132 & 49,878 & 0 \\
134 & 62,372 & 0 \\
136 & 79,362 & 6 \\
138 & 98,541 & 0 \\
140 & 121,354 & 0 \\
142 & 151,201 & 0 \\
144 & 186,611 & 15 \\
146 & 225,245 & 0 \\
148 & 277,930 & 0 \\
150 & 335,569 & 0 \\
152 & 404,667 & 17 \\
154 & 489,646 & 0 \\
156 & 586,264 & 0 \\
158 & 697,720 & 0 \\
160 & 836,497 & 40 \\
162 & 989,495 & 0 \\
164 & 1,170,157 & 1 \\
166 & 1,382,953 & 0 \\
168 & 1,628,029 & 45 \\
170 & 1,902,265 & 1 \\
172 & 2,234,133 & 0 \\
174 & 2,601,868 & 0 \\
176 & 3,024,383 & 89 \\
178 & 3,516,365 & 0 \\
180 & 4,071,832 & 1 \\
182 & 4,690,880 & 0 \\
184 & 5,424,777 & 117 \\
186 & 6,229,550 & 0 \\
188 & 7,144,091 & 0 \\
190 & 8,187,581 & 0 \\
192 & 9,364,975 & 199 \\
194 & 10,659,863 & 0 \\
196 & 12,163,298 & 0 \\
198 & 13,809,901 & 0 \\
200 & 15,655,672 & 271 \\
\midrule Total & 120,622,224 & 816\\\bottomrule
\end{longtable}
% END GENERATED COUNTS

\subsection{The 816 classes and their relation to A3}
The saved taxonomy U assigns disjoint explanations:
\begin{center}\small
\begin{tabular}{lr}\toprule
Explanation & Classes\\\midrule
$\GC(2,0)$ images &812\\
$\GC(3,0)[C_{20}]$, at $C_{180}$ &1\\
$C1(2,3)=F(2,3)$, at $C_{164}$ &1\\
$O2(2,3)=G(2,2,3)$, at $C_{170}$ &1\\
$C1(3,2)=F(3,2)$, at $C_{184}$ &1\\
Unexplained in this finite census &0\\\bottomrule
\end{tabular}\end{center}
The GC2 comparison includes 116 at $C_{184}$, 199 at $C_{192}$, and
271 at $C_{200}$. The three noninteger-GC examples were identified by
full canonical-code equality with their constructions. Here
``noninteger GC'' means not $\GC(d,0)[F]$ for an ordinary fullerene
$F$ and $d\ge2$, not the complement of all GC images. The C1 and O2
identifications are with existing catalogue families, not new family
constructions. Their general absence arguments are separate from the
individual certificates and this finite accounting.

In the theorem's interior domain, $20T\le200$ permits only $(2,1)$,
$T=7$, $C_{140}$, which is IPR and positive, consistently with $m=1$.
All 121,354 IPR $C_{140}$ outputs are positive. The boundary points
$C_{60}=\GC(1,1)[C_{20}]$, $C_{80}=\GC(2,0)[C_{20}]$, and
$C_{180}=\GC(3,0)[C_{20}]$ are not interior-theorem cases. The census
therefore does not verify a broad range of interior A3 parameters.

\subsection{Independent comparison and the local data package}
Section 10.3 fixes the independent receipt boundary at $C_{182}$,
with updated cumulative totals withheld pending the coordinator's
final check. The published-data plan adopted for local preparation
contains an English README and schema, a 71-row CSV, an 816-row JSONL
index, 816 certificate JSON files, audit summaries, decision and
checking code with dependency hashes, and generator metadata.
The certificate files total 60,758,514 uncompressed bytes. Generator
source and binaries, raw output streams, and all positive graphs are
excluded. Package-relative references retain both original and derived
hashes. Local archive hashes and the exact validation scope are recorded
separately; no archive is uploaded by this preparation. Licensing,
third-party reproduction, author details, and a publication venue remain
release decisions.

\clearpage
\section{The negative B orientation: coordinates, cuts, and ownership}\label{app:B}
\claimid{app:B}
This appendix expands Section 6 without new parameter conditions. It
covers every piece of every fine triangle, not just vertices. Assign
piece interiors to parent faces and use the agreement of neighbouring
affine gluing maps on boundaries. The facial-triangle prerequisite,
primitive parent edges, and the exclusion of earlier parent vertices
are as in the body and Section 7.

\subsection{Two coordinate systems}
Let $(x,y)$ be B's original fine coordinates and $(\alpha,\beta)$ its
original parent coordinates. Its triangle sequence lies in
$0\le x\le b$, $0\le y\le1$, with
$T\alpha=(k+l)x+ly$, $T\beta=-lx+ky$.
For the two slanted fan types the actual faces are rotated by $R^{-1}$.
For calculation only, rotate back to common coordinates:
\begin{equation}\label{eq:B1}
\begin{gathered}
(X,Y)=(-y,x+y),\qquad(u,v)=(-\beta,\alpha+\beta),\\
Tu=lx-ky,\quad Tv=kx+(k+l)y,\quad
T(u+v)=(k+l)x+ly.
\end{gathered}\tag{B1}
\end{equation}
Thus $-1\le X\le0$, $0\le X+Y\le b$, not A's rectangle.
The original parent-face names are recovered by
\begin{equation}\label{eq:B2}
R^{-1}U(i,j)=D(i+j,-i-1),\quad
R^{-1}D(i,j)=U(i+j+1,-i-1).\tag{B2}
\end{equation}
This planar coordinate change does not rotate quotient labels; use
the actual B gluing tables for each face map.

\subsection{Outer fan: \texorpdfstring{$k=ql+r$}{k=ql+r}}
Here retain $(\alpha,\beta)$, with strip $-1\le\beta\le0$.
Call the six positive expressions in Section 6.1 $F_1,\ldots,F_6$.
Throughout the rectangle,
\[
\beta\ge-lb/T>-2,\quad \beta\le k/T<1,\quad
\alpha+\beta\le(kb+k+l)/T<q+2.
\]
The first bound uses $2T-lb=F_1+k+l$ and the last uses $F_3$.
If $\beta>0$, then $x<ky/l$, $\alpha\ge0$, and
$\alpha+\beta=[kx+(k+l)y]/T<y/l\le1/l<1$, placing the initial
piece in $U(0,0)$. At $\beta=0$, $\alpha=y/l$. At $\beta=-1$,
\[
 \alpha=q+1+(r+y)/l\in(q+1,q+2),
\]
using $r\ge2$, $r+1\le l-1$; the cut beyond $(q+2,-1)$ is avoided.
If $\beta<-1$, then $\alpha+\beta=(y-k\beta)/l>k/l>q$.
The resulting region $-2<\beta<-1$, $q<\alpha+\beta<q+2$ is
covered by
\[
D(q+1,-2),\ U(q+2,-2),\ D(q+2,-2),\ U(q+3,-2).
\]
The line $\alpha+\beta=q+1$ splits the first pair from the second;
$\alpha=q+2$ or $q+3$ identifies the individual triangle. The first
crossing is the common edge of $U(q+1,-1)$ and $D(q+1,-2)$.
At $(b,1)$,
\[
T(-1-\beta)=F_6>0,\qquad
T(\alpha+\beta-q-1)=F_5+k+l>0.
\]
Thus $H_B$ applies on the last two faces. Its translation differs from
the B translation in (7) by $-(5,0)$. The early point $(a,0)$ has
$\beta=-la/T$ and $T-la=F_5+k+l>0$, so it is inside the old belt.
Together with ordering and Section 7, this gives $2b+1\to2a$.

\subsection{Positive-side fans: \texorpdfstring{$k=cl+w$, $l=sw+d$}{k=cl+w, l=sw+d}}
Let $c=1,2$, with $s\ge2,1$, $2\le d\le w-2$, and
$W_v=cs+c+1$. The common normal is
$n=v-cu=\alpha+(c+1)\beta$, with
\[
Tn=wx+[(c+1)k+l]y\ge0,\qquad -y=wu-ln.
\]
There is no initial negative-normal B piece. At $n=1$,
$u=s+(d-y)/w\in(s,s+1)$. For $n>1$,
$u>(l-1)/w>s$, so exit is on the cell-$s$ side of the cut from $W$.
Call the quantities (8) $P_1,\ldots,P_9$. Their exact uses are:
\begin{center}\small
\begin{tabular}{ll}\toprule
Required bound & Positive numerator\\\midrule
$v<W_v+1$ & $P_4$\\
$u+v<s+W_v+2$ & $P_5+(k-l)$\\
$u<s+2$ & $P_9+k+l$\\
$v>W_v$ at $(b,0)$ & $P_6+l$\\
$u+v>s+W_v+1$ at $(b,0)$ & $P_7$\\
$n<1$ at $(a,1)$ & $P_8$\\
$0<a<b<T$ & $P_1,P_2,P_3$\\\bottomrule
\end{tabular}\end{center}
For complete ownership, partition the $n>1$ portion into three cases.
If $u\le s+1,v\le W_v$, then $s<u\le s+1$ and
$cu+1<v\le cs+c+1$, within the original full cell triangles
$U(s,j),D(s,j)$, $j=cs,\ldots,cs+c$. Here ``cell'' includes the full
glued triangles, not just their intersection with the slanted strip.
If $u\le s+1,v>W_v$, the bound $v<W_v+1$ gives $U(s,W_v)$ or
$D(s,W_v)$, divided by $u+v=s+W_v+1$.
If $u>s+1$, then $v>W_v$ and $u+v<s+W_v+2$, giving $U(s+1,W_v)$.
This covers every piece. The last point belongs to the latter two faces.
In B's original coordinates the departure and added faces are
\[
\begin{gathered}
U(s+W_v,-s-1);\\
D(s+W_v,-s-1),\ U(s+W_v+1,-s-1),\ D(s+W_v+1,-s-2).
\end{gathered}
\]
Conjugating the common map gives
$H_B(z)=R^{-1}z+(s+1,W_v)$. The difference from the translation
$(s-5,cs+5)$ in (9) is $(-6,4-c)$: minus one B period,
$(6,-3)$ or $(6,-2)$. At $(a,1)$, $0<n<1$, so the early point
uses the same belt map. This establishes $2b\to2a+1$ without an
appeal to analogy with A.

\subsection{Negative-side fan: \texorpdfstring{$k=2l-w$, $l=sw+d$}{k=2l-w, l=sw+d}}
Here $s\ge2$, $2\le d\le w-2$, and
\[
n=v-2u=\alpha+3\beta,\quad Tn=-wx+(3k+l)y,\quad -y=-wu-ln.
\]
At $n=-1$, $u=s+(d+y)/w\in(s,s+1)$, using $d+1\le w-1$.
For $n<-1$, $u>(l-1)/w>s$, before the cut from $(s+1,2s+1)$.
Writing (10) as $M_1,\ldots,M_8$, the whole rectangle has
$u<s+2$ from $M_4+k+l$ and $u+v<3s+3$ from $M_5+(k-l)$.
At $(b,1)$ the lower bounds are $M_6$ for $u>s+1$ and $M_7+l$
for $u+v>3s+2$. At $(a,0)$, $T-wa=M_8+2k+3l>0$ gives $n>-1$;
$M_1,M_2,M_3$ order the indices.

Since $v=(ku+y)/l$, $u>s$ implies $v>2s-1+d/l>2s-1$.
For $u>s+1$,
\[
v>k(s+1)/l=2s+[(s-1)w+2d]/l>2s.
\]
If $n<-1,u\le s+1$, then $2s-1<v<2u-1\le2s+1$, within
$U(s,j),D(s,j)$ for $j=2s-1,\ldots,2s+1$. If $u>s+1$, the
bounds put the point in $U(s+1,2s),D(s+1,2s),U(s+1,2s+1)$:
below $v=2s+1$, split at $u+v=3s+2$; above it, use the last $U$.
The last point is in the final two faces. Under (B2), these faces are
\[
\begin{gathered}
 D(s,2s)\mapsto U(3s+1,-s-1),\\
 U(s+1,2s)\mapsto D(3s+1,-s-2),\\
 D(s+1,2s)\mapsto U(3s+2,-s-2),\\
 U(s+1,2s+1)\mapsto D(3s+2,-s-2).
\end{gathered}
\]
For the initial $n>0$ piece, return to original coordinates. From
$y=ln-w\beta$,
\[
0\le\alpha=\frac{3y-(3l-w)n}{w}<3/w\le3/4,\qquad
0\le\alpha+\beta=\frac{2y-kn}{w}<2/w\le1/2.
\]
It lies in $U(0,0)\cup D(0,-1)$, with $D$ central and $U$ additional.
Assign $-1\le n\le0$ to the central belt. This covers the rectangle.
The original map is $H_B(z)=Rz+(2s+1,-3s-2)$. The translation
in (11) differs by $(-6,2)=-R^{-1}(2,4)$. The early point has
$-1<n=-wa/T<0$, proving $2b+1\to2a$ at the same ordinary vertex.

\begin{samepage}\section{Terminal positive caps by point reflection}\label{app:C}
\claimid{app:C}
\subsection{Action on the sequence and parent triangles}
For $W=(A,B)$ let $J_W(z)=W-z$. In fine coordinates this is
$(1,Q)-(x,y)$ for A, $(Q,1)-(x,y)$ for B, and $(Q-1,1)-(x,y)$
for AR B. It preserves the relevant rectangle and gives
$P_{E-j}=P_E-P_j$. Consecutive triangles and their cut pieces map to
the terminal triangles in reverse order. Direct substitution gives\end{samepage}
\begin{equation}\label{eq:C1}
J_WU(i,j)=D(A-i-1,B-j-1),\quad
J_WD(i,j)=U(A-i-1,B-j-1).\tag{C1}
\end{equation}
This is planar geometry, not quotient reflection symmetry. Terminal
face maps and labels are taken from the fixed cap table.

\subsection{AR B, \texorpdfstring{$k=ql-1$}{k=ql-1}}
Here $W=(q+1,-1)$, $n=v$, and $n(J_Wz)=-1-n(z)$.
The initial $n>0$ piece has $u\ge0$, $v>0$,
$u+v\le y/l\le1/l<1$, so lies in $U(0,0)$.
For the terminal $n<-1$ piece, $z'=J_Wz$ satisfies
\[
u'=q+1-u\ge0,\quad v'=-1-v>0,\quad
u'+v'=q-u-v\le1/l<1.
\]
Thus the terminal cap is $D(q,-2)$, while the initial central
$D(0,-1)$ maps to central $U(q,-1)$. Their shared edge is
$[(q,-1),(q+1,-1)]$. Reflection preserves the central normal interval.
The initial U phase zero and terminal D phase $q-1\bmod5$ have labels
distinct from the central belt and each other. Width one was handled
earlier; choose the plus branch at width two. Containing each cut piece
does not require all vertices of a triangle touching that piece to
lie in the same cap face.

\subsection{Second A and diagonal A, \texorpdfstring{$l=sw-1$}{l=sw-1}}
For $c=2$, $w=k-2l$; for $c=1$, $w=k-l$. Set $n=v-cu$ and
$W=(s,cs+1)$. The minus branch uses $w\ge3$, and
$n(J_Wz)=1-n(z)$. On the initial $n<0$ piece,
\[
0\le u<1/w,\qquad -u\le v<cu,\qquad u+v<(c+1)/w\le1.
\]
It lies in central $U(0,0)$ and additional $D(0,-1)$. Reflecting a
terminal $n>1$ point gives
\[
\begin{gathered}
0\le s-u<1/w,\quad -(s-u)\le cs+1-v<c(s-u),\\
(c+1)s+1-u-v<(c+1)/w\le1.
\end{gathered}
\]
The terminal central and additional faces are $D(s-1,cs)$ and
$U(s-1,cs+1)$, sharing $[(s-1,cs+1),(s,cs+1)]$.
The extra-face phases are $s-1\bmod2$ for second A and $s\bmod3$
for diagonal A. The fixed table exhausts them. Rectangle invariance
and piecewise reflection cover every terminal triangle piece.

\subsection{Folded second B, \texorpdfstring{$l=sw-1$, $w=2l-k\ge3$}{l=sw-1, w=2l-k >= 3}}
Here $W=(3s-1,-s)$, $n=u+3v$, and $n(J_Wz)=-1-n(z)$.
For the initial $n>0$ piece with $0\le y\le1$,
\[
0\le u=\frac{3y-(3l-w)n}{w}<3/w\le1,\quad
0\le u+v=\frac{2y-kn}{w}<2/w<1.
\]
Thus $U(0,0)$ is additional and $D(0,-1)$ central. Reflecting the
terminal $n<-1$ piece gives
\[
0\le3s-1-u<3/w,\quad0\le2s-1-u-v<2/w,\quad n(z')>0.
\]
The terminal additional face is $D(3s-2,-s-1)$ and the central face
$U(3s-2,-s)$, sharing $[(3s-2,-s),(3s-1,-s)]$.
The two phases $s\bmod2$ occur in the table. The central normal
interval is preserved and boundary pieces glue along the shared edge.

\subsection{From containment to legality on the surface}
This appendix proves cap containment and gluing as a prose argument.
Distinct face labels still depend on the finite table; exclusion of
periodic identifications on Section 5.5; and degrees on $Q<T$ and the
candidate congruences. It does not by itself complete review of the
classification. The Japanese v4 B and C prose received the coordinator's
AI review in \record{49430e1c}; this English transcription and the
selected-formula machine checks do not constitute human peer review,
a second independent general proof, or full formalization.

\section*{Formal lemma identifiers}
These identifiers refer to the exact types and evidence in \record{A3-v6.2-machine/claims.json}.
The 41 A3 arithmetic declarations, the conditional finite-path theorem,
and three selected prototype declarations have distinct entries.
\par\noindent\label{lean:unit_endpoint}\record{unit_endpoint}
\claimid{lean:unit_endpoint}
\par\noindent\label{lean:second_plus}\record{second_plus}
\claimid{lean:second_plus}
\par\noindent\label{lean:second_minus}\record{second_minus}
\claimid{lean:second_minus}
\par\noindent\label{lean:diagonal_plus}\record{diagonal_plus}
\claimid{lean:diagonal_plus}
\par\noindent\label{lean:diagonal_minus}\record{diagonal_minus}
\claimid{lean:diagonal_minus}
\par\noindent\label{lean:folded_second_plus}\record{folded_second_plus}
\claimid{lean:folded_second_plus}
\par\noindent\label{lean:folded_second_minus}\record{folded_second_minus}
\claimid{lean:folded_second_minus}
\par\noindent\label{lean:outer_collisions}\record{outer_collisions}
\claimid{lean:outer_collisions}
\par\noindent\label{lean:positive_fan_collisions}\record{positive_fan_collisions}
\claimid{lean:positive_fan_collisions}
\par\noindent\label{lean:negative_fan_collisions}\record{negative_fan_collisions}
\claimid{lean:negative_fan_collisions}
\par\noindent\label{lean:no_small_multiple}\record{no_small_multiple}
\claimid{lean:no_small_multiple}
\par\noindent\label{lean:outer_u_bound}\record{outer_u_bound}
\claimid{lean:outer_u_bound}
\par\noindent\label{lean:outer_v_bound}\record{outer_v_bound}
\claimid{lean:outer_v_bound}
\par\noindent\label{lean:candidate_order}\record{candidate_order}
\claimid{lean:candidate_order}
\par\noindent\label{lean:outer_endpoint_order}\record{outer_endpoint_order}
\claimid{lean:outer_endpoint_order}
\par\noindent\label{lean:au_c1_a_pos}\record{au_c1_a_pos}
\claimid{lean:au_c1_a_pos}
\par\noindent\label{lean:au_c1_b_minus_a_pos}\record{au_c1_b_minus_a_pos}
\claimid{lean:au_c1_b_minus_a_pos}
\par\noindent\label{lean:au_c1_N_minus_b_pos}\record{au_c1_N_minus_b_pos}
\claimid{lean:au_c1_N_minus_b_pos}
\par\noindent\label{lean:au_c1_v_upper_pos}\record{au_c1_v_upper_pos}
\claimid{lean:au_c1_v_upper_pos}
\par\noindent\label{lean:au_c1_sum_upper_pos}\record{au_c1_sum_upper_pos}
\claimid{lean:au_c1_sum_upper_pos}
\par\noindent\label{lean:au_c1_v_late_lower_pos}\record{au_c1_v_late_lower_pos}
\claimid{lean:au_c1_v_late_lower_pos}
\par\noindent\label{lean:au_c1_sum_late_lower_pos}\record{au_c1_sum_late_lower_pos}
\claimid{lean:au_c1_sum_late_lower_pos}
\par\noindent\label{lean:au_c1_early_inside_pos}\record{au_c1_early_inside_pos}
\claimid{lean:au_c1_early_inside_pos}
\par\noindent\label{lean:au_c1_endpoint_bound_pos}\record{au_c1_endpoint_bound_pos}
\claimid{lean:au_c1_endpoint_bound_pos}
\par\noindent\label{lean:au_c2_a_pos}\record{au_c2_a_pos}
\claimid{lean:au_c2_a_pos}
\par\noindent\label{lean:au_c2_b_minus_a_pos}\record{au_c2_b_minus_a_pos}
\claimid{lean:au_c2_b_minus_a_pos}
\par\noindent\label{lean:au_c2_N_minus_b_pos}\record{au_c2_N_minus_b_pos}
\claimid{lean:au_c2_N_minus_b_pos}
\par\noindent\label{lean:au_c2_v_upper_pos}\record{au_c2_v_upper_pos}
\claimid{lean:au_c2_v_upper_pos}
\par\noindent\label{lean:au_c2_sum_upper_pos}\record{au_c2_sum_upper_pos}
\claimid{lean:au_c2_sum_upper_pos}
\par\noindent\label{lean:au_c2_v_late_lower_pos}\record{au_c2_v_late_lower_pos}
\claimid{lean:au_c2_v_late_lower_pos}
\par\noindent\label{lean:au_c2_sum_late_lower_pos}\record{au_c2_sum_late_lower_pos}
\claimid{lean:au_c2_sum_late_lower_pos}
\par\noindent\label{lean:au_c2_early_inside_pos}\record{au_c2_early_inside_pos}
\claimid{lean:au_c2_early_inside_pos}
\par\noindent\label{lean:au_c2_endpoint_bound_pos}\record{au_c2_endpoint_bound_pos}
\claimid{lean:au_c2_endpoint_bound_pos}
\par\noindent\label{lean:au_neg_a_pos}\record{au_neg_a_pos}
\claimid{lean:au_neg_a_pos}
\par\noindent\label{lean:au_neg_b_minus_a_pos}\record{au_neg_b_minus_a_pos}
\claimid{lean:au_neg_b_minus_a_pos}
\par\noindent\label{lean:au_neg_N_minus_b_pos}\record{au_neg_N_minus_b_pos}
\claimid{lean:au_neg_N_minus_b_pos}
\par\noindent\label{lean:au_neg_u_upper_pos}\record{au_neg_u_upper_pos}
\claimid{lean:au_neg_u_upper_pos}
\par\noindent\label{lean:au_neg_sum_upper_pos}\record{au_neg_sum_upper_pos}
\claimid{lean:au_neg_sum_upper_pos}
\par\noindent\label{lean:au_neg_u_late_lower_pos}\record{au_neg_u_late_lower_pos}
\claimid{lean:au_neg_u_late_lower_pos}
\par\noindent\label{lean:au_neg_sum_late_lower_pos}\record{au_neg_sum_late_lower_pos}
\claimid{lean:au_neg_sum_late_lower_pos}
\par\noindent\label{lean:au_neg_early_inside_pos}\record{au_neg_early_inside_pos}
\claimid{lean:au_neg_early_inside_pos}
\par\noindent\label{lean:finite_literal_path_decision}\record{finite_literal_path_decision}
\claimid{lean:finite_literal_path_decision}

\par\noindent\label{lean:unit_width_path}\record{unit_width_path}
\claimid{lean:unit_width_path}
\par\noindent\label{lean:chartAdj_iff}\record{chartAdj_iff}
\claimid{lean:chartAdj_iff}
\par\noindent\label{lean:faces_connected}\record{faces_connected}
\claimid{lean:faces_connected}

\section*{Acknowledgements and AI-use statement (draft)}
The human author should select and approve this statement before any
release. Proposed wording: ``The author used Codex and Claude Code
to assist with literature organization, mathematical exploration,
proof drafting, implementation, finite verification, formalization of
specified lemmas, and English exposition. The author is responsible
for the published account, its provenance, and the stated limits of
verification. AI systems are not listed as authors. AI prose reviews
are identified separately from human review. Only the explicitly
listed Lean propositions are described as formally proved; finite
checks and independently constructed tests have their stated domains.
No claim of complete human checking of every argument is made.''
This paragraph is a proposal, not a statement already approved for publication.

\section*{Reproducibility records}
Relative repository paths are archival identifiers, not claims of
public availability. The earlier proof drafts are
\record{A3-unit-residue-proof-AR.md}, \record{A3-qge3-proof-AT.md},
\record{A3-q2-positive-proof-AT.md}, \record{A3-complete-proof-AU.md},
and \record{A3-endpoint-order-AU.md}. Finite geometry is preserved in
\record{review-AW/manifest.json}, its cap table, and the three atlases
named in Section 6. Review execution hashes are in
\record{AX-validation.json}, \record{A3-v2-validation.json},
\record{A3-v3-validation.json}, and \record{A3-v4-validation.json}.
The English correspondence and validation record accompany this source.
The machine-readable claim layer distinguishes dependencies and
evidence statuses; its identifiers match the labels above. Hashes
establish file identity, not the truth of mathematical contents.

\begin{thebibliography}{10}
\bibitem{BBS} A. Bille, V. Buchstaber, and E. Spodarev,
\emph{Some Open Mathematical Problems on Fullerenes}, arXiv version
2024; published 2025. \href{https://doi.org/10.1021/acs.jcim.4c01997}{doi:10.1021/acs.jcim.4c01997}.
\bibitem{DD} M. Dutour and M. Deza,
\emph{Goldberg--Coxeter Construction for 3- and 4-valent Plane Graphs},
2004. \href{https://doi.org/10.37236/1773}{doi:10.37236/1773}.
\bibitem{ZZ} M. Deza and M. Dutour,
\emph{Zigzag Structures of Simple Two-faced Polyhedra},
arXiv:math/0212352v2, 2003.
\bibitem{Growing} G. Brinkmann, D. Franceus, P. W. Fowler, and J. E. Graver,
\emph{Growing fullerenes from seed}, 2006.
\href{https://doi.org/10.1016/j.cplett.2006.07.040}{doi:10.1016/j.cplett.2006.07.040}.
\bibitem{Patches} G. Brinkmann, P. W. Fowler, and C. Justus,
\emph{A Catalogue of Isomerization Transformations of Fullerene Polyhedra},
2003. \href{https://doi.org/10.1021/ci020069l}{doi:10.1021/ci020069l}.
\bibitem{Capsids} A. Luque and D. Reguera,
\emph{The Structure of Elongated Viral Capsids}, 2010.
\href{https://doi.org/10.1016/j.bpj.2010.02.051}{doi:10.1016/j.bpj.2010.02.051}.
\bibitem{Pentagonal} F. Sadre-Marandi and P. Das,
\emph{Extension of Caspar--Klug theory to higher order pentagonal polyhedra},
2018. \href{https://doi.org/10.1515/cmb-2018-0001}{doi:10.1515/cmb-2018-0001}.
\bibitem{Knots} M. Deza, M. Dutour, and P. W. Fowler,
\emph{Zigzags, Railroads, and Knots in Fullerenes}, 2004.
\href{https://doi.org/10.1021/ci049955h}{doi:10.1021/ci049955h}.
\bibitem{BGW} T. Bogart, J. Goodrick, and K. Woods,
\emph{Parametric Presburger arithmetic: logic, combinatorics, and
quasi-polynomial behavior}, 2017, Theorem 1.15.
\end{thebibliography}
\end{document}
